\documentclass[12pt]{article}

% PAGE FORMAT
\oddsidemargin 0.0in
\evensidemargin 0.0in
\textwidth 6.5in
\textheight 8.5in
\topmargin -0.25in

% FOR TOGGLING ON CHOICES OF WHAT TO PRINT
\def \showcodes{}     % Comment out to NOT see codes
%\def \showspace{}     % Comment out to NOT have space after problems.
%\def \showsolutions{} % Comment out to NOT see solutions
%\def \showrubric{}     % Comment out to NOT see grading rubric
\def \showtags{} % Comment out to NOT see tags
%%%%
% NOTE: showsolutions negates effects of: showrubric and showspace
% NOTE: showrubric negates effects of: showspace


% KEY HEADER FILE
% *****************************************************************
% INCLUDE THE FOLLOWING IN ALL CLASS WORK, HOME WORK, TESTS
% *****************************************************************

% PACKAGES
\usepackage{amsmath}
\usepackage{latexsym}
\usepackage{amssymb}
\usepackage{ifthen}
\usepackage{multicol}
\usepackage{enumitem}
\usepackage{graphicx}
\usepackage{float}
\usepackage{tikz}
\usetikzlibrary{calc}
\usetikzlibrary{positioning}
\usepackage{pgfplots}
%\usepackage{C:/Users/Kerry/Dropbox/MY_STUFF/MY_MATERIALS/Programs/Python/ExerciseManual/ExerciseData/cancel}
\usepackage{cancel}

% COMMANDS
%\renewcommand{\labelenumi}{\alph{enumi})} % To make enumerates use letters
\newcommand{\im}{\ensuremath{\mathbf{i}}} % For imaginary ``i''
\newcommand{\code}[1]{ {\underline{\textsc(#1)}} } % For coding problems
\newcommand{\ans}[2]{{\bf Answer: } \framebox(#1,#2){} } % For making answer box

%For commenting out
\newcommand{\junk}[1]{}



% Command for inserting solution figure
\newcommand{\solfig}[1]{
\begin{figure}[H] 
    \centering
    \includegraphics[scale=.7]{C:/Users/Kerry/Dropbox/MY_STUFF/MY_MATERIALS/Programs/Python/ExerciseManual/ExerciseData/FigsSolutions/SOLUTIONxxx#1.pdf} 
\end{figure}
}


% KEY EXERCISE COMMAND: 6 ARGUMENTS -
% #1: Code
% #2: Amount of space for student work
% #3: The actual exercise text
% #4: Solution (I expect this to usually link to PDF figure)
% #5: Grading Rubric (I expect this to usually link to PDF figure which is mofication of solution)
% #6: Database Tags
\newcommand{\ex}[6]
{
\begin{minipage}{\textwidth}
\item
	\ifdefined \showcodes      % Display code (#1) if variable is on.
		#1 \\ 
	\fi
	\ifdefined \showtags        % Display database tags (#6) if variable is on.
	    #6 \\
   \fi
	#3                                  % Always show the problem (#3)
	\ifdefined \showsolutions   % Show solution if variable (#4) is on

		\vspace{2mm}
		SOLUTION.
		\vspace{2mm}
		#4
		\vspace{2mm}
	\else
		\ifdefined \showrubric    % Show rubric if variable (#5) on and NOT showing solutions
			#5 
		\else
			\ifdefined \showspace  % Give space if variable (#2) on and NOT showing solutions, nor rubric
				\vspace{#2 mm} 
			\fi
		\fi
	\fi
	\vspace{2mm}                   % Always have ending delimiter
	\hrule
	\vspace{2mm}
\end{minipage}
}



% Or put hard link to your file something like in the following.
%\input{C:/Users/[NAME]/ ... /ExerciseData/HeaderExerciseFileStandard.tex}   



\begin{document}

\noindent
{\textbf{Assignment \# .....}}\\

\noindent
{\bf Instructions:} .....




\begin{enumerate}

 \ex
{\code{DISCx25x08x15x009}}{10}
{ 
 How many rows appear in a truth table for each of these
compound propositions?

\begin{multicols}{2}
\begin{enumerate}

\item $X \lor \neg X$

\item $A \rightarrow (A \land \neg B)$

\item $(P \lor \neg Q) \leftrightarrow ( (P \land Q) \rightarrow \neg P)$

\item $(A \lor B) \rightarrow (A \lor B \land \neg C)$

\end{enumerate}
\end{multicols}


}{

 


\begin{align*}
\text{1. For } X \lor \neg X:& \quad \text{1 variable } (X) \Rightarrow 2^1 = 2 \text{ rows} \\
\text{2. For } A \rightarrow (A \land \neg B):& \quad \text{2 variables } (A, B) \Rightarrow 2^2 = 4 \text{ rows} \\
\text{3. For } (P \lor \neg Q) \leftrightarrow ( (P \land Q) \rightarrow \neg P):& \quad \text{2 variables } (P, Q) \Rightarrow 2^2 = 4 \text{ rows} \\
\text{4. For } (A \lor B) \rightarrow (A \lor B \land \neg C):& \quad \text{3 variables } (A, B, C) \Rightarrow 2^3 = 8 \text{ rows}
\end{align*}

}{}{topics: DISC PROPLOGIC}


 \ex
{\code{DISCx25x08x15x010}}{20}
{ 
 Construct a truth table for the following proposition:
$X \lor \neg X$


}{

 
$
\begin{array}{|c|c|c|}
\hline
X & \neg X & X \lor \neg X \\
\hline
\text{T} & \text{F} & \text{T} \\
\text{F} & \text{T} & \text{T} \\
\hline
\end{array}
$
}{}{topics: DISC PROPLOGIC}



 \ex
{\code{DISCx25x09x11x001}}{30}
{ 
 Construct a truth table for the following proposition and state whether it is a tautology, contradiction, or contingency:
$(X \land \neg Y) \rightarrow (\neg Y)$


}{

 

\[
\begin{array}{c|c|c|c|c}
X & Y & \neg Y & X \land \neg Y & (X \land \neg Y) \rightarrow \neg Y \\
\hline
\text{T} & \text{T} & \text{F} & \text{F} & \text{T} \\
\text{T} & \text{F} & \text{T} & \text{T} & \text{T} \\
\text{F} & \text{T} & \text{F} & \text{F} & \text{T} \\
\text{F} & \text{F} & \text{T} & \text{F} & \text{T} \\
\end{array}
\]

$(X \land \neg Y) \rightarrow (\neg Y)$ is a tautology.
}{}{type: TEST , topics: DISC PROPLOGIC}




 \ex
{\code{QUANTx25x08x18x007}}{20}
{  Determine the truth value of each of these statements if
the domain consists of all rational numbers.  In
both cases, what if we changed the domain to be just the integers?

\begin{multicols}{2} 
\begin{enumerate}

\item $\forall x \exists y \ x \cdot y = 1$

\item  $\forall x (x > 0 \rightarrow (x - 3) > 0 )$



\end{enumerate}
\end{multicols} 

}{

 

1. For the domain of rational numbers:

- Statement: $\forall x \exists y \ x \cdot y = 1$

- Truth value: True, since for any rational number $x$, there exists a rational $y = \frac{1}{x}$ such that $x \cdot y = 1$.

2. For the domain of integers:

- Statement: $\forall x \exists y \ x \cdot y = 1$

- Truth value: False, since for $x = 0$, there is no integer $y$ such that $0 \cdot y = 1$.

3. For the domain of rational numbers and integers:

- Statement: $\forall x (x > 0 \rightarrow (x - 3) > 0 )$

- Truth value: False, since when $x = 1$, $x > 0$ but $(x - 3) = -2 \not> 0$.
}{}{topics: DISC QUANT}


 \ex
{\code{QUANTx26x02x04x001}}{20}
{  Determine the truth value of each of these statements if
the domain consists of all integers

\begin{multicols}{2} 
\begin{enumerate}


\item $\exists u \forall x \  x < \sqrt{u}$

\item $\forall x \exists u  \   x < \sqrt{u} $


\end{enumerate}
\end{multicols} 

}{

 

$\exists u \forall x \  x < \sqrt{u}$ is false.

$\forall x \exists u \  x < \sqrt{u}$ is true.
}{}{type: HW , topics: DISC QUANT}





\end{enumerate} 
\end{document}